% Part: incompleteness
% Chapter: representability-in-q
% Section: c-representable

\documentclass[../../../include/open-logic-section]{subfiles}

\begin{document}

\olfileid{inc}{req}{cre}
\olsection{د $C$ تابعې په $\Th{Q}$ کښې تمثيلېدونکې دي}

بايد وښيو چې په $C$ کښې هره تابعه په $\Th{Q}$ کښې تمثيلېدونکې ده۔ په پاے کښې دا ښودل پکار دي چې د $C$ هرې $k$-ځايه تابعې $f(x_0,\dots,x_{k-1})$ ته داسې فارمول $!A_f(x_0,\dots,x_{k-1},y)$ څنګه ګومارلے شو چې هغه تمثيل کړي۔

\begin{lem}
که طبيعي عددونه $n$ او $m$ راکړل شوي وي او $n \neq m$ وي، نو $Q \Proves \num n \neq \num m$۔
\end{lem}

\begin{proof}
پر $n$ استقرا وکاروئ، څو وښيئ چې د هر $m$ دپاره، که $n \neq m$ وي، نو $Q \Proves \eq/[\num n][\num m]$۔

په بنسټيز حالت کښې $n = 0$ دے۔ که $m$ له $0$ سره برابر نۀ وي، نو د کوم طبيعي عدد $k$ دپاره $m = k + 1$ دے۔ يو اصل لرو چې وايي $\lforall[x][\eq/[0][x']]$۔ د کمیت ټاکونکي د يو اصل په مرسته، د $x$ پر ځاے $\num k$ په اېښودلو سره، $\eq/[0][\num k']$ ترلاسه کولے شو۔ خو $\num k'$ هماغه $\num m$ دے۔

په استقرايي ګام کښې فرض کولے شو چې دعوه د $n$ دپاره صحيح ده، او بيا $n+1$ په پام کښې نيسو۔ $m$ دې هر طبيعي عدد وي۔ دوه امکانونه شته: يا $m = 0$ دے، يا د کوم $k$ دپاره $m = k+1$ دے۔ لومړے حالت د پورته طريقې په شان حل کېږي۔ په دويم حالت کښې فرض کړئ چې $n+1 \neq k+1$ دے۔ نو $n \neq k$ دے۔ د $n$ دپاره د استقرايي فرض له مخې $\Th{Q} \Proves \eq/[\num n][\num k]$ لرو۔ يو اصل وايي $\lforall[x][\lforall[y][\eq[x'][y'] \lif \eq[x][y]]]$۔ د کمیت ټاکونکي د يو اصل په کارولو $\eq[\num n'][\num k'] \lif \eq[\num n][\num k]$ لرو۔ د قضيو د منطق په مرسته په $\Th{Q}$ کښې $\eq/[\num n][\num k] \lif \eq/[\num n'][\num k']$ ترلاسه کولے شو۔ د مودوس پوننس په کارولو $\eq/[\num n'][\num k']$ ترلاسه کېږي، چې همدا مو مطلوبه پايله ده، ځکه $\num k'$ هماغه $\num m$ دے۔
\end{proof}

پام وکړئ چې دا مرستندويه قضيه ډېره لويه دعوه نۀ کوي: په اصل کښې وايي چې $\Th{Q}$ ثابتولے شي چې بېلې عددنښې بېل څيزونه ښيي۔ مثلاً $\Th{Q}$ ثابتوي چې $0'' \neq 0'''$۔ خو د دې عمومي ثبوت يو څه احتياط غواړي۔ دا هم په پام کښې وساتئ چې که څه هم استقرا کاروو، دا استقرا د $\Th{Q}$ \emph{څخه بهر} ده۔

صفر، تالي، جمع، ضرب، د عينيت مشخصه تابعه او پروجکشنونه تمثيلولے شو۔ په هر حالت کښې مناسب تمثيلوونکے فارمول ډېر نېغ په نېغه دے؛ مثلاً صفر په دې فارمول تمثيلېږي:
\[
y = 0,
\]
تالي په دې فارمول تمثيلېږي:
\[
x_0' = y,
\]
او جمع په دې فارمول تمثيلېږي:
\[
x_0 + x_1 = y.
\]
کار دا دے چې وښيو $\Th{Q}$ اړوند بيانونه ثابتولے شي؛ مثلاً دا ويل چې جمع په پورته فارمول تمثيلېږي، دا ښودل غواړي چې د طبيعي عددونو د هرې جوړې $m$ او $n$ دپاره $\Th{Q}$ دا ثابتوي:
\[
\eq[\num n + \num m][\num {n+m}]
\]
او
\[
\lforall[y][(\eq[(\num n + \num m)][y] \lif \eq[y][\num {n+m}])].
\]

د ترکيب په اړه څۀ ويلے شو؟ فرض کړئ چې $h$ داسې تعريف شوے دے:
\[
h(x_0,\dots,x_{l-1}) = f(g_0(x_0,\dots,x_{l-1}), \dots,
g_{k-1}(x_0,\dots,x_{l-1})).
\]
او هغه فارمولونه $!A_f, !A_{g_0}, \dots, !A_{g_{k-1}}$ مو له مخکې موندلي دي چې په ترتيب سره تابعې $f,g_0,\dots,g_{k-1}$ تمثيلوي۔ بيا د $h$ د تمثيل دپاره $!A_h$ تعريفولے شو، په دې توګه چې $!A_h(x_0,\dots,x_{l-1},y)$ دا فارمول وټاکو:
\begin{multline*}
  \lexists[z_0\dots][\lexists[z_{k-1}][(!A_{g_0}(x_0,\dots,x_{l-1},z_0) \land
  \dots \land !A_{g_{k-1}}(x_0,\dots,x_{l-1},z_{k-1}) \land]]\\
  !A_f(z_0,\dots,z_{k-1},y)).
\end{multline*}

په پاے کښې بې‌حده پلټنه په پام کښې نيسو۔ فرض کړئ چې $g(x,\vec z)$ منظمه ده او په $\Th{Q}$ کښې، مثلاً په فارمول $!A_g(x,\vec z,y)$، تمثيلېدونکې ده۔ $f$ دې په $f(\vec z) = \mu x \; g(x,\vec z)$ تعريف شي۔ غواړو داسې فارمول $!A_f(\vec z,y)$ ومومو چې $f$ تمثيل کړي۔ يو طبيعي انتخاب دا دے:
\[
!A_f(\vec z,y) \ident !A_g(y,\vec z,0) \land \lforall[w][(w < y
\lif \lnot !A_g(w,\vec z,0))].
\]

د ځينو نورو مرستندويه قضيو په کارولو ښودل کېدے شي چې دا انتخاب کار کوي؛ مثلاً:

\begin{lem}
د هر متغير $x$ او هر طبيعي عدد $n$ دپاره $\Th{Q}$ دا ثابتوي: $\eq[(x'
  + \num n)][(x + \num n)']$۔
\end{lem}

بيا هم دا يادونه پکار ده چې دا تر هغې دعوې کمزورې خبره ده چې $\Th{Q}$، $\lforall[x][\lforall[y][(x' + y) = (x +y)']]$ ثابتوي؛ د اثبات دغه دعوه ناسمه ده، نۀ په معياري طبيعي عددونو کښې خپله حسابي مساوات۔

\begin{proof}
ثبوت د معمول په شان پر $n$ په استقرا دے۔ په بنسټيز حالت کښې $n = 0$ دے؛ بايد وښيو چې $\Th{Q}$، $\eq[(x'+0)][(x+0)']$ ثابتوي۔ خو دا لرو:
\begin{eqnarray*}
& & \eq[(x' + 0)][x'] \quad \text{له اصل ۴ څخه} \\
& & \eq[(x + 0)][x] \quad \text{له اصل ۴ څخه} \\
& & \eq[(x+0)'][x'] \quad \text{له کرښې ۲ څخه} \\
& & \eq[(x' + 0)][(x + 0)'] \quad \text{له کرښو ۱ او ۳ څخه}
\end{eqnarray*}
په استقرايي ګام کښې فرض کولے شو چې $\eq[(x' + \num n)][(x + \num n)']$ مو په $\Th{Q}$ کښې استنتاج کړے دے۔ ځکه چې $\num{n+1}$ هماغه $\num n'$ دے، بايد وښيو چې $\Th{Q}$، $\eq[(x' + \num n')][(x + \num n')']$ ثابتوي۔ د مساواتو دا لړۍ په $\Th{Q}$ کښې استنتاج کېدے شي:
\begin{eqnarray*}
(x' + \num n') & \eq & (x' + \num n)' \quad \text{اصل ۵} \\
& \eq & (x + \num n')' \quad \text{له استقرايي فرض او اصل ۵ څخه}
\end{eqnarray*}
\end{proof}

\begin{lem}
\begin{enumerate}
\item $\Th{Q}$، $\lnot (x < \num 0)$ ثابتوي۔
\item د هر طبيعي عدد $n$ دپاره $\Th{Q}$ دا ثابتوي:
\[
x < \num {n+1} \lif (\eq[x][\num 0] \lor \dots \lor \eq[x][\num n]).
\]
\end{enumerate}
\end{lem}

\begin{proof}
لومړۍ برخه او د دويمې برخې يوه برخه په غير رسمي ډول کوو، يعنې يوازې دا اشارې ورکوو چې رسمي !!{derivation} څنګه جوړ شي۔

د لومړۍ برخې دپاره، د $<$ د تعريف له مخې، بايد په $\Th{Q}$ کښې $\lnot \lexists[y][\eq[(y'+x)][0]]$ ثابت کړو؛ دا د لومړۍ درجې منطق د اصولو او قواعدو په کارولو له $\lforall[y][\eq/[(y' + x)][0]]$ سره همارزښته ده۔ مفکوره دا ده: فرض کړئ چې $\eq[(y'+x)][0]$ دے۔ که $x$، 0 وي، $\eq[(y' + 0)][0]$ لرو۔ خو د $\Th{Q}$ د اصل ۴ له مخې $\eq[(y' + 0)][y']$ لرو، او د اصل ۲ له مخې $\eq/[y'][0]$ لرو، چې تناقض دے۔ نو $\lforall[y][\eq/[(y' +x)][0]]$۔ که $x$، $0$ نۀ وي، د اصل ۳ له مخې داسې $z$ شته چې $\eq[x][z']$ دے۔ خو بيا $\eq[(y' + z')][0]$ لرو۔ د اصل~۵ له مخې $\eq[(y' + z)'][0]$ لرو، چې بيا هم د اصل~۲ سره په تناقض کښې دے۔

د دويمې برخې دپاره پر $n$ استقرا وکاروئ۔ بنسټيز حالت په پام کښې نيسو، چې $n =0$ دے۔ په دې حالت کښې بايد $x < \num 1 \lif \eq[x][\num 0]$ وښيو۔ فرض کړئ چې $x < \num 1$ دے۔ نو د $<$ د تعريفي اصل له مخې $\lexists[y][\eq[(y'+x)][0']]$ لرو۔ فرض کړئ چې $y$ دغه خاصيت لري؛ نو $\eq[y'+x][0']$ لرو۔

بايد $\eq[x][0]$ وښيو۔ د اصل ۳ له مخې، که $x$، 0 نۀ وي، د کوم $z$ دپاره له $z'$ سره برابر دے۔ بيا $\eq[(y' + z')][0']$ لرو۔ د $\Th{Q}$ د اصل~۵ له مخې $\eq[(y'+z)'][0']$ لرو۔ د اصل ۱ له مخې $\eq[(y' + z)][0]$ لرو۔ خو دا د تعريف له مخې $z < 0$ معنا لري، چې له لومړۍ برخې سره په تناقض کښې دے۔
\end{proof}

\begin{explain}
ښودلې مو ده چې د محاسبه کېدونکو تابعو سټ د هغو تابعو د سټ په توګه ځانګړے کېدے شي چې په $\Th{Q}$ کښې تمثيلېدونکې دي۔ په حقيقت کښې ثبوت تر دې عام دے۔ د تمثيلېدنې له تعريف څخه په اسانۍ ليدل کېږي چې هره تيوري چې $\Th{Q}$ غځوي (يا $\Th{Q}$ پکې تفسيرېدے شي)، محاسبه کېدونکې تابعې تمثيلولے شي؛ خو برعکس، په هر هغه !!{derivation} نظام کښې چې د اثبات مفهوم محاسبه کېدونکے وي، هره تمثيلېدونکې تابعه محاسبه کېدونکې ده۔ نو مثلاً د محاسبه کېدونکو تابعو سټ د هغو تابعو د سټ په توګه ځانګړے کېدے شي چې د پيانو په حساب، يا ان د زرمیلو--فرېنکل په سټ تيورۍ کښې تمثيلېږي۔ لکه ګوډل چې يادونه کړې، دا يو څه هېښوونکې ده۔ وبه وينو چې د ثبوت‌وړتيا پوښتنې د ټاکلې تيورۍ پر وړاندې ډېرې حساسې دي؛ په ټوليز ډول، هر څومره چې اصول پياوړي وي، هغومره ډېر څه ثابتولے شئ۔ خو د اصولي تيوريو په يوه پراخه ډله کښې تمثيلېدونکې تابعې کټ مټ محاسبه کېدونکې تابعې دي۔ \emph{OLCMP-099: د برعکس جهت دپاره د تيورۍ سازګاري او د بېلابېلو عددنښو د نابرابرۍ اثبات هم لازم دي۔ يوازې د اثبات محاسبه‌وړتيا بس نۀ ده؛ په ناسازګار کلاسيک نظام کښې د هرې تابعې، حتٰی د نامحاسبه کېدونکې تابعې، د تمثيل دواړه شرطونه له تناقض څخه ثابتېږي۔ د قوي تيوريو وروستۍ پرتله هم د سازګارۍ په فرض ده۔}
\end{explain}

\end{document}
